\documentclass[../../main.tex]{subfiles}
\begin{document}

{\bf [解]}
$a>0$とする．$A(a,0)$, $P=(x,y)\neq A$とおくと，
\begin{align}
 \overrightarrow{AP}=
  \begin{pmatrix}
   x-a \\ y
  \end{pmatrix} \label{eq:1}
\end{align}
である．半直線$AP$上の点$Q$が$\dfrac{AQ}{AP}\le 2$を満たす時，
\begin{align}
 0 \le t \le 2 \label{eq:2}
\end{align}
なる実数$t$に対して
\begin{align}
 \overrightarrow{AQ} = t \overrightarrow{AP} = 
  t\begin{pmatrix}
   x-a \\ y
  \end{pmatrix} \\
\therefore
 Q(a+t(x-a),ty) \label{eq:3}
\end{align}
とおける．

題意の条件は，この範囲のすべての$t$に対して
\begin{align}
\dfrac{QP}{OQ}\le\dfrac{AP}{OA}\quad(OQ\neq 0) \label{eq:4}
\end{align}
が成り立つことである．$QP=|1-t|\cdot AP$，$OA=a$だから，\cref{eq:4}に代入して（$AP>0$より）
\begin{align}
 a|1-t|\le OQ=\sqrt{(a+t(x-a))^2+(ty)^2} \quad (\because \text{\cref{eq:3}})
\end{align}
である．

両辺とも非負なので2乗して整理すると
\begin{align}
 & a^2(1-t)^2\le(a+t(x-a))^2+t^2y^2 \\
 & t\bigl[(x^2-2ax+y^2)t+2ax\bigr]\ge 0 \label{eq:4} 
\end{align}
である．

$t=0$のとき\cref{eq:4}は自明に成立するから，$t\in(0,2]$で$h(t):=(x^2-2ax+y^2)t+2ax\ge0$が成り立てばよい．$h$は$t$の1次以下の式（線形）だから，区間$[0,2]$上で$h\ge0$となるための必要十分条件は両端点で$h\ge0$となることである．従って条件は
\begin{align}
 \begin{cases}
 h(0)=2ax\ge 0 \\
 h(2)=2(x^2-ax+y^2)\ge 0
 \end{cases} \\
 \begin{cases}
 h(0)=x\ge 0 \quad (\because a>0)\\
 h(2)=\Bigl(x-\dfrac{a}{2}\Bigr)^2+y^2\ge \Bigl(\dfrac{a}{2}\Bigr)^2 \ge 0
 \end{cases} \label{eq:5}
\end{align}
である．

最後に$P\neq A$および$OQ\neq0$つまり$P\neq O$だから，これら$2$点を除外すればよく，求める領域$D$は
\begin{align}
D:\ \Bigl\{(x,y)\ \Big|\ x\ge0,\ \Bigl(x-\dfrac{a}{2}\Bigr)^2+y^2\ge\Bigl(\dfrac{a}{2}\Bigr)^2\Bigr\}\setminus\{(0,0),\,(a,0)\}
\end{align}
すなわち，$O(0,0)$と$A(a,0)$を直径の両端とする円の外部（境界は含むが，$O$，$A$の2点は除く）．
これを図示すると下図斜線部（ただし白丸は除く）である．

\begin{figure}[htb]
\begin{center}
\begin{tikzpicture}[scale=1.6, >=stealth]

  % --- 1. パラメータ設定 (例: a = 2.4, 半径 R = a/2 = 1.2) ---
  \def\aVal{2.4}
  \pgfmathsetmacro{\rVal}{\aVal/2}

  % 描画・ハッチング範囲
  \def\xMin{0}
  \def\xMax{4.0}
  \def\yRange{2.1}

  % --- 2. 領域 D の斜線ハッチング (x >= 0 から円の内部をくり抜く) ---
  \begin{scope}
    % x >= 0 の半平面をクリップ
    \clip (\xMin, -\yRange) rectangle (\xMax, \yRange);
    % 全体をハッチング
    \fill[pattern=north east lines, pattern color=blue!60] 
      (\xMin, -\yRange) rectangle (\xMax, \yRange);
    % 円の内部 (x - a/2)^2 + y^2 < (a/2)^2 を白くマスクしてくり抜く
    \fill[white] (\rVal, 0) circle (\rVal);
  \end{scope}

  % --- 3. 座標軸 ---
  \draw[->, thick] (-0.8, 0) -- (\xMax + 0.5, 0) node[right] {$x$};
  \draw[->, thick] (0, -\yRange - 0.2) -- (0, \yRange + 0.3) node[above] {$y$};
  \node[below left, font=\small] at (0,0) {$O$};

  % --- 4. 境界線の描画 (実線: 境界を含む) ---
  % y 軸 (x = 0)
  \draw[ultra thick, blue!80!black] (0, -\yRange) -- (0, \yRange);

  % 円 (x - a/2)^2 + y^2 = (a/2)^2
  \draw[ultra thick, blue!80!black] (\rVal, 0) circle (\rVal);

  % --- 5. 中心・目盛り・ラベル ---
  % 円の中心 (a/2, 0)
  \fill[black] (\rVal, 0) circle (1.0pt) node[below, font=\small] {$\dfrac{a}{2}$};

  % 点 A(a, 0)
  \node[below right, font=\small] at (\aVal, 0) {$A(a,0)$};

  % 境界円の方程式ラベル
  \node[blue!90!black, font=\scriptsize, align=center] at (\rVal, 0.45) 
    {$\Bigl(x-\dfrac{a}{2}\Bigr)^2+y^2=\Bigl(\dfrac{a}{2}\Bigr)^2$};

  % 領域ラベル
  \node[blue!90!black, font=\Large] at (3.0, 1.2) {$D$};

  % --- 6. 除外点 O, A の白丸プロット（境界は含むがこの2点のみ除く） ---
  \filldraw[fill=white, draw=black, thick] (0, 0) circle (2.2pt);
  \filldraw[fill=white, draw=black, thick] (\aVal, 0) circle (2.2pt);

\end{tikzpicture}
\end{center}
\caption{点 $P$ の存在領域 $D$（斜線部．境界を含み，点 $O(0,0)$ と $A(a,0)$ のみ除く）}
\end{figure}

\end{document}
